\section{Deducción del gradiente de política}

\subsection{Función objetivo y factorización de la trayectoria}

La función objetivo que maximiza el gradiente de política es la esperanza de la recompensa de la trayectoria:
$$J(\theta) = \mathbb{E}_{\tau \sim p_\theta(\tau)} \left[ \sum_t r(s_t, a_t) \right],$$
donde $p_\theta(\tau) = p(s_1)\prod_t \pi_\theta(a_t\mid s_t)p(s_{t+1}\mid s_t,a_t)$. Consideramos el parámetro de política $\theta$ como el único objeto de optimización; los demás términos solo determinan la distribución de muestreo.

\begin{knowledgebox}{Sentido de la factorización de la trayectoria}
\begin{itemize}
    \item La parte dinámica $p(s_{t+1}\mid s_t,a_t)$ no depende de $\theta$, por lo que no afecta al gradiente;
    \item la política $\pi_\theta$ determina cómo se muestrea, y $J(\theta)$ también es una esperanza sobre la política de muestreo;
    \item la recompensa se acumula a lo largo de toda la trayectoria y forma una señal global, que hay que descomponer en forma diferencial.
\end{itemize}
\end{knowledgebox}

\subsection{Log-Gradient Trick}

Usando la identidad $\nabla_\theta p_\theta(\tau) = p_\theta(\tau) \nabla_\theta \log p_\theta(\tau)$, se puede escribir:
$$\nabla_\theta J(\theta) = \mathbb{E}_{\tau \sim p_\theta(\tau)} \left[ \nabla_\theta \log p_\theta(\tau) \cdot r(\tau) \right].$$
Al desarrollar $\log p_\theta(\tau)$ se obtiene: $\log p(s_1) + \sum_t \left[\log \pi_\theta(a_t\mid s_t) + \log p(s_{t+1}\mid s_t,a_t)\right]$,
donde solo $\log \pi_\theta$ depende de los parámetros, con lo cual se obtiene:
$$\nabla_\theta J(\theta) = \mathbb{E}_{\tau \sim p_\theta(\tau)} \left[ \left(\sum_t \nabla_\theta \log \pi_\theta(a_t\mid s_t)\right) \cdot r(\tau) \right].$$

\begin{importantbox}{No hace falta conocer el modelo del entorno}
La dinámica del entorno $p(s_{t+1}\mid s_t,a_t)$ se cancela mediante el log-gradient trick, de modo que podemos muestrear trayectorias directamente y estimar el objetivo en un entorno model-free.
\end{importantbox}

\subsection{Aproximación de Monte Carlo: REINFORCE}

En la práctica, se reemplaza la esperanza por una estimación de Monte Carlo:
$$\nabla_\theta J(\theta) \approx \frac{1}{N} \sum_{i=1}^{N} \left(\sum_t \nabla_\theta \log \pi_\theta(a_{i,t}\mid s_{i,t})\right) \cdot r(\tau_i).$$
Esta fórmula es precisamente el estimador de gradiente de REINFORCE. El flujo del algoritmo es: (1) muestrear $N$ trayectorias; (2) calcular la recompensa de cada trayectoria; (3) actualizar los parámetros con la estimación anterior.

\begin{table}[H]
\centering
\small
\begin{tabular}{p{0.3\textwidth}p{0.45\textwidth}p{0.2\textwidth}}
\toprule
Símbolo & Significado & Observación \\
\midrule
$\sum_t r(s_t,a_t)$ & Recompensa total de la trayectoria & Supervisión a nivel de trayectoria \\
$\nabla_\theta \log \pi_\theta$ & Gradiente de la log-verosimilitud de la política & Se puede diferenciar automáticamente \\
$N$ & Número de trayectorias muestreadas & Afecta la varianza y el cómputo \\
\bottomrule
\end{tabular}
\caption{Cantidades clave del estimador REINFORCE}
\end{table}

\subsection{Resumen de la sección}

El núcleo del gradiente de política consiste en escribir la esperanza de la recompensa de la trayectoria como el producto de la log-verosimilitud por la recompensa, de modo que el gradiente se pueda estimar directamente sin modelar la dinámica. REINFORCE ofrece la implementación de Monte Carlo más elemental, pero con una varianza considerable; las secciones siguientes discuten formas de mejorarla.
